
\subsubsection*{Settle carried liabilities in the period that pays them}
The new opening can now support ordinary events and tax payments. Its two
liabilities still refer to earlier assessments, but the cash-flow statement
must record only payments made in the new financial period. Reusing the old
settlement list as current cash movement would count earlier payments again.
The continuation therefore keeps each historical vintage unchanged and records
new settlements against its separate current-period opening.

Suppose the company pays 3 on 1 January 2027 and explicitly allocates 1 to the
liability opening at $15/4$ and 2 to the liability opening at 5. Their remaining
amounts become $11/4$ and 3. Total payable falls from $35/4$ to $23/4$, while
cash falls from $715/2$ to $709/2$. The entry debits tax payable and credits cash
by 3; it creates no current-tax expense. The new period's operating cash flow
is $-3$. Historical payments remain available for inspection without entering
that sum.

Each payment names the stable liability identities and supplies positive exact
amounts whose sum equals the payment. Each allocation must fit that liability's
remaining current-period amount, and the total must fit available cash. The
caller also supplies affirmative payment authority. No oldest-first allocation,
proration, borrowing or reassessment is inferred. A second payment can exhaust
the remaining $11/4$ and 3, leaving both liability records at zero and cash at
$1395/4$. Keeping the zero records preserves which assessments were settled.

Ordinary events act on the same full financial state. If depreciation of 2
follows the first payment, the right-of-use asset falls from 78 to 76 and
recorded profit becomes $-2$. Assets are then $1351/2$, liabilities $443/4$ and
equity $2259/4$. The indirect operating cash calculation adds back depreciation
and includes the reduction in opening tax payable, giving $-3$, the same amount
as the dated cash payment. A retained inventory lot can also be sold through
the selected return method; its history and tax adjustments remain attached to
the resulting book.

Operations must have unused lifetime identities and chronological dates inside
the new financial interval. A rejected allocation, missing authority or cash
shortage stops that operation and preserves every preceding successful event,
together with the exact unprocessed sequence. Applying a proposed result
recomputes and compares the whole result, including its opening, financial
history, liability records and recognition archive. A nonportable request,
such as one containing a floating-point amount, is rejected before execution.

These events do not reassess current or deferred tax. The carried loss
quantities, recognition findings and historical measurements remain intact,
while both new-period assessments stay unassessed. A reusable annual design
must next explain how a fresh assessment joins all outstanding liabilities and
how another boundary preserves their complete histories.
